%============================================================
% BAB 1: PENDAHULUAN
% Struktur (Proposal dan Buku PA):
%   1.1 Latar Belakang
%   1.2 Permasalahan
%   1.3 Batasan Masalah
%   1.4 Tujuan
%   1.5 Manfaat
%   1.6 Sistematika Penulisan
%============================================================
\chapter{PENDAHULUAN}

\section{LATAR BELAKANG}

Perkembangan arsitektur perangkat lunak modern yang beralih dari aplikasi monolitik ke sistem berbasis \textit{microservice} dan \textit{container} membuat proses migrasi maupun pengoperasiannya menjadi kompleks dan menuntut sumber daya yang signifikan \cite{trabelsi2025systematic}. Lingkungan produksi kontemporer dapat terdiri dari banyak layanan yang saling bergantung dan menghasilkan volume data telemetri yang besar, sehingga menuntut pendekatan \textit{monitoring} yang lebih cerdas dibandingkan metode konvensional yang masih bergantung pada intervensi manusia secara manual \cite{hrusto2025monitoring}.

Tantangan ini juga dihadapi oleh infrastruktur \textit{backend} aplikasi FoodLAB di Politeknik Elektronika Negeri Surabaya (PENS), aplikasi pemesanan makanan daring untuk civitas akademika kampus. \textit{Observability}, yakni kemampuan sistem untuk memperlihatkan kondisi internalnya melalui data telemetri berupa \textit{metrics}, \textit{log}, dan \textit{traces}, menjadi kualitas kritis yang harus dimiliki infrastruktur \textit{backend} FoodLAB \cite{trinity2023}. Kondisi ini diperberat oleh proses \textit{rewrite} arsitektur \textit{backend} FoodLAB dari Laravel dan MySQL menuju ASP.NET Core dan PostgreSQL yang sedang berlangsung, serta akumulasi \textit{code smells} dan \textit{technical debt} akibat pengembangan estafet oleh berbagai generasi mahasiswa \cite{yulanda2025refactoring}. Selama periode transisi tersebut, kondisi operasional sistem menjadi sulit diprediksi sehingga mekanisme deteksi anomali yang adaptif menjadi lebih penting dibandingkan pendekatan berbasis \textit{threshold} statis.

Pendekatan \textit{monitoring} DevOps tradisional yang mengandalkan \textit{threshold} statis dan respons insiden manual memiliki keterbatasan signifikan. Pada infrastruktur produksi skala menengah dapat muncul ribuan metrik dengan puluhan hingga ratusan jenis \textit{alert}, sehingga administrator rentan mengalami \textit{alert fatigue} dan mulai mengabaikan \textit{alert} yang sebenarnya penting \cite{voutsas2024mitigating}. Di sisi lain, mekanisme \textit{self-healing} konvensional yang berbasis \textit{runbook} dan skrip otomasi hanya efektif untuk skenario kegagalan yang telah diantisipasi sebelumnya, dan tidak mampu bernalar untuk menangani pola kegagalan baru maupun kegagalan \textit{cascading} yang kompleks \cite{liang2023self}. Keterbatasan ini menjadi semakin kritis pada aplikasi seperti FoodLAB yang terus berkembang dan ditangani oleh tim pengembang yang berganti secara periodik.

Kemajuan \textit{Large Language Models} (LLM) dan \textit{framework agentic AI} membuka peluang pendekatan \textit{monitoring} yang lebih adaptif. LLM terbukti mampu melakukan penalaran kontekstual untuk tugas yang kompleks dan tidak terstruktur \cite{xu2025large}, sementara integrasi kecerdasan adaptif berbasis AI ke dalam alur \textit{DevSecOps} mampu meningkatkan kemampuan deteksi dan respons secara signifikan \cite{dinh2025enhancing}. Arsitektur \textit{multi-agent} berlapis yang terdiri dari \textit{Perception Layer}, \textit{Reasoning Layer}, dan \textit{Action Layer} juga terbukti efektif untuk dekomposisi tugas \textit{root cause analysis} \cite{fu2025leveraging}. Berdasarkan hal tersebut, penelitian ini mengusulkan \textit{framework} \textit{Agentic DevOps Monitoring Agents} (ADMA) untuk mengumpulkan dan menganalisis telemetri multi-sinyal FoodLAB secara \textit{real-time}, melakukan diagnosis akar masalah otomatis berbasis LLM, serta mengeksekusi remediasi secara otonom, sehingga \textit{observability} dan kelangsungan pengembangan \textit{backend} FoodLAB dapat lebih terjaga.

\section{RUMUSAN MASALAH}

Infrastruktur \textit{backend} aplikasi FoodLAB di Politeknik Elektronika Negeri Surabaya (PENS) menghadapi beberapa permasalahan terkait \textit{monitoring} dan pemulihan insiden yang berpotensi menghambat ketersediaan layanan dan pengembangan jangka panjang:

\begin{enumerate}[label=\arabic*., leftmargin=*]
  \item \textbf{\textit{Monitoring} yang Reaktif dan Bergantung pada Manusia}\\
        Sistem \textit{monitoring} konvensional yang digunakan pada infrastruktur FoodLAB masih berbasis \textit{threshold} statis yang menghasilkan \textit{false-positive rate} tinggi dan \textit{alert fatigue}. Operator harus melakukan investigasi manual terhadap berbagai sumber data seperti \textit{log}, metrik, dan \textit{resource} secara terpisah, yang menyulitkan tim dalam mendeteksi dan memulihkan insiden secara cepat. Kondisi ini berdampak pada memanjangnya dua metrik kunci manajemen insiden: \textit{Mean Time to Detect} (MTTD) dan \textit{Mean Time to Recovery} (MTTR). MTTD adalah rata-rata waktu yang dibutuhkan untuk mendeteksi suatu masalah sejak masalah tersebut terjadi, sedangkan MTTR adalah rata-rata waktu yang dibutuhkan untuk memulihkan sistem ke kondisi normal setelah insiden terdeteksi. Keduanya terdampak secara signifikan ketika visibilitas sistem tidak memadai \cite{voutsas2024mitigating}.

  \item \textbf{Ketiadaan \textit{Root Cause Analysis} Otomatis}\\
        Ketika terjadi anomali pada \textit{backend} FoodLAB, seperti \textit{container} \textit{crash}, \textit{query} \textit{database} yang lambat, atau \textit{resource exhaustion}, operator harus secara manual menelusuri \textit{log} dan metrik dari berbagai komponen untuk mengidentifikasi akar masalah. Proses ini rawan kesalahan dan memakan waktu lama, terutama pada kegagalan \textit{cascading} yang melibatkan beberapa komponen sekaligus \cite{trinity2023}.

  \item \textbf{Tidak Efisiennya \textit{Root Cause Analysis} Akibat Ketiadaan Penyaringan Sinyal Awal}\\
        Proses \textit{root cause analysis} yang memanggil \textit{Large Language Model} secara langsung tanpa mekanisme penyaringan awal berisiko menimbulkan biaya operasional yang tinggi dan waktu respons yang lambat ketika sistem menangani banyak sinyal anomali secara bersamaan. Tanpa mekanisme untuk membatasi kandidat investigasi yang diteruskan ke LLM, seluruh daftar \textit{tool} yang tersedia akan dikonsumsi setiap siklus penalaran, sehingga konsumsi \textit{token} tumbuh secara tidak terkendali seiring pertumbuhan skala sistem \cite{voutsas2024mitigating}.

  \item \textbf{Ketiadaan Mekanisme Perbaikan Kode Otomatis}\\
        Ketika akar masalah teridentifikasi berasal dari cacat pada kode sumber aplikasi, belum ada mekanisme otomatis yang menghasilkan usulan perbaikan kode. Proses perbaikan sepenuhnya bergantung pada intervensi manual pengembang, termasuk untuk kasus perbaikan sederhana yang berulang. Kondisi ini diperparah oleh \textit{knowledge gap} akibat pergantian tim pengembang secara berkala, yang memperpanjang waktu pemulihan bahkan untuk masalah yang sebenarnya dapat diotomasi \cite{liang2023self}.

  \item \textbf{Tidak Adanya Mekanisme \textit{Self-Healing}}\\
        \textit{Backend} FoodLAB belum memiliki kemampuan untuk secara otomatis melakukan tindakan korektif saat terjadi insiden. Setiap kegagalan, mulai dari \textit{container} yang \textit{crash} hingga konfigurasi yang salah, memerlukan intervensi manual dari tim pengembang. Hal ini menjadi semakin problematik mengingat pengembangan FoodLAB dilakukan secara estafet oleh tim yang berganti setiap periode, sehingga \textit{knowledge gap} mengenai infrastruktur dapat memperlambat proses pemulihan \cite{liang2023self}.
\end{enumerate}

Untuk mengatasi permasalahan tersebut, diperlukan suatu sistem \textit{monitoring} cerdas yang mampu melakukan deteksi anomali secara proaktif, analisis \textit{root cause} secara otomatis, dan eksekusi remediasi secara otonom. \textit{Framework} \textit{Agentic DevOps Monitoring Agents} (ADMA) diusulkan sebagai solusi yang mengintegrasikan kemampuan AI berbasis LLM untuk mengubah paradigma \textit{monitoring} dari reaktif menjadi proaktif dan otonom pada infrastruktur \textit{backend} FoodLAB.

\section{BATASAN MASALAH}

Dalam pelaksanaan Proyek Akhir ini, terdapat beberapa batasan yang ditetapkan agar penelitian lebih terfokus dan terarah, yaitu:

\begin{enumerate}[label=\arabic*., leftmargin=*]
  \item Pemantauan \textit{database} oleh \textit{Metric Perception Agent} (MPA) hanya mencakup PostgreSQL. \textit{Backend} FoodLAB versi lama yang masih berjalan di atas MySQL (Laravel) tidak termasuk dalam cakupan pemantauan MPA, meskipun \textit{Log Perception Agent} dan \textit{System Resource Agent} tetap dapat memantau \textit{container} dan \textit{resource host} pada kedua \textit{environment} secara \textit{database-agnostic}.
  \item Evaluasi sistem dilakukan pada lingkungan \textit{staging} yang mereplikasi infrastruktur FoodLAB generasi baru berbasis ASP.NET Core dan PostgreSQL, bukan pada infrastruktur produksi FoodLAB yang melayani pengguna aktif.
  \item Infrastruktur target yang dipantau berbasis Docker mandiri, tidak mencakup lingkungan \textit{container orchestration} seperti Kubernetes.
  \item \textit{Reasoning Engine} dievaluasi menggunakan Claude API sebagai jalur utama. Penggunaan model LLM lokal (Qwen3-14B) dibahas sebagai alternatif desain untuk skenario keterbatasan \textit{resource} GPU, bukan bagian dari evaluasi eksperimen utama pada Bab IV.
  \item \textit{Action Executor} dibatasi pada tindakan \textit{restart container}, perbaikan konfigurasi, \textit{kill query database}, dan pembuatan \textit{merge request} GitLab; tidak mencakup tindakan skala infrastruktur lain seperti \textit{auto-scaling} atau migrasi \textit{server}.
  \item \textit{Code Fix Agent} hanya aktif untuk \textit{container} berkategori \textit{application} yang memiliki kode sumber di GitLab, dan tidak aktif untuk \textit{container} berkategori \textit{database} atau \textit{infrastructure}.
  \item Data operasional yang dikirim ke layanan LLM pihak ketiga dibatasi pada \textit{log}, \textit{query}, dan cuplikan kode teknis yang relevan dengan insiden yang sedang diinvestigasi, bukan data pribadi pengguna aplikasi FoodLAB.
  \item Penelitian dilaksanakan pada periode Juli hingga Agustus 2026 sesuai jadwal Proyek Akhir.
\end{enumerate}

\section{TUJUAN}

Tujuan dari proyek ini adalah untuk meningkatkan \textit{observability} dan kemampuan \textit{self-healing} pada infrastruktur \textit{backend} aplikasi FoodLAB melalui penerapan \textit{framework} \textit{Agentic DevOps Monitoring Agents} (ADMA). Secara spesifik, proyek ini bertujuan untuk:

\begin{enumerate}[label=\arabic*., leftmargin=*]
  \item Merancang dan mengimplementasikan sistem \textit{multi-agent} otonom yang mampu melakukan \textit{monitoring} \textit{real-time} pada infrastruktur \textit{backend} FoodLAB, mencakup PostgreSQL, Docker \textit{container} \textit{logs}, dan \textit{resource} sistem.
  \item Mengembangkan mekanisme \textit{reasoning} berbasis LLM yang mampu melakukan \textit{root cause analysis} secara otomatis berdasarkan data telemetri multi-sinyal yang dikumpulkan oleh \textit{Perception Agents}.
  \item Merancang mekanisme korelasi sinyal (\textit{Alert Correlator}) dan penyaringan kandidat investigasi (\textit{Pre-Filter}) berbasis aturan deterministik untuk mengurangi beban pemanggilan LLM, serta menghitung tingkat keyakinan diagnosis (\texttt{ValidityScore}) sebagai dasar keputusan apakah \textit{pipeline} investigasi perlu diulang sebelum rencana remediasi diusulkan kepada operator.
  \item Mengimplementasikan \textit{Code Fix Agent} yang mampu menghasilkan usulan perbaikan kode secara otomatis melalui pembuatan \textit{merge request} pada \textit{repository} GitLab.
  \item Mengimplementasikan kemampuan \textit{self-healing} yang memungkinkan sistem secara otonom mengeksekusi tindakan remediasi seperti \textit{restart container} dan perbaikan konfigurasi.
  \item Menyediakan \textit{control plane} terpusat melalui \textit{backend} C\# ASP.NET Core sebagai sarana pendukung tercapainya tujuan 1 sampai 5, yaitu untuk pengelolaan \textit{agent}, visualisasi status sistem, dan \textit{audit trail} seluruh aktivitas \textit{monitoring} dan remediasi.
\end{enumerate}

\section{MANFAAT}

Proyek ini diharapkan memberikan manfaat signifikan baik dari sisi teoritis maupun praktis dalam pengelolaan infrastruktur dan keberlanjutan pengembangan aplikasi FoodLAB di PENS.

\textbf{Manfaat Teoritis}

Dari sisi akademik, penelitian ini memberikan kontribusi berupa implementasi konkret dari konsep \textit{Agentic AI} untuk \textit{DevOps monitoring} pada lingkungan produksi nyata. Pendekatan ini mendemonstrasikan bahwa paradigma \textit{monitoring} otonom berbasis \textit{multi-agent} dan LLM dapat diterapkan pada skala aplikasi kampus dengan infrastruktur berbasis Docker, tidak hanya terbatas pada lingkungan Kubernetes skala \textit{enterprise}. Penelitian ini juga berkontribusi pada pengembangan metodologi evaluasi sistem \textit{agentic} yang terukur, mencakup metrik efisiensi token, akurasi \textit{root cause analysis}, dan laju pemulihan otonom.

\textbf{Manfaat Praktis}

Dari sisi operasional, penerapan ADMA diharapkan dapat mengurangi MTTD dan MTTR secara signifikan melalui deteksi anomali yang proaktif dan remediasi yang otonom. Sistem ini memungkinkan \textit{monitoring} 24/7 tanpa bergantung pada ketersediaan operator manusia, sehingga meningkatkan \textit{availability} layanan FoodLAB bagi civitas akademika PENS.

Dari sisi pengembangan, ADMA menyediakan mekanisme \textit{observability} yang komprehensif melalui \textit{dashboard} dan \textit{audit trail} yang terpusat. Tim pengembang baru yang mengambil alih proyek FoodLAB dapat dengan mudah memahami kondisi infrastruktur, melihat riwayat insiden dan remediasi, serta mendapat manfaat dari kemampuan \textit{self-healing} yang mengurangi beban pengelolaan manual.

Secara keseluruhan, manfaat utama dari proyek ini adalah terciptanya infrastruktur \textit{backend} FoodLAB yang lebih \textit{resilient}, \textit{observable}, dan \textit{self-healing}, sehingga mendukung kelangsungan pengembangan aplikasi secara berkelanjutan.

\section{SISTEMATIKA PENULISAN}

Proposal proyek akhir ini disusun dalam tiga bab utama yang masing-masing membahas aspek penting dari perencanaan penelitian dan pengembangan sistem ADMA pada infrastruktur \textit{backend} aplikasi FoodLAB. Sistematika penulisan proposal ini adalah sebagai berikut:

\input{\SistematikaFile}
